% FILE: 02_lineage_and_context.tex
% Project: ggen — Governance Generation Engine

\section{Lineage and Context}
\label{sec:ggen-lineage}

\subsection{2016 Language-Model Lesson}
\label{sec:ggen-lineage-2016}

The research program's 2016 language-model experiment established that local text imitation does
not conserve consequence: a system that reproduces the surface form of an utterance does not
thereby reproduce the causal structure that gave rise to that utterance. Applied to capital-market
artifacts, this lesson is sharp. A system that generates acquisition deck slides by imitating
prior deck language may reproduce the persuasive register of board communication without
reproducing the evidentiary grounding that legitimizes the claims. The result is an artifact that
looks board-admissible but is not independently verifiable.

ggen is the engineering response to this lesson. Rather than generating artifact text by
imitation, ggen manufactures artifacts by \emph{derivation}: each output is a deterministic
function of SPARQL query results applied to a template with strict variable mode enabled
(\texttt{strict\_variables = true} in \texttt{ggen.toml}). A slide asserting fitness $\geq 0.95$
is not an imitation of a claim about fitness; it is the rendered output of a SPARQL
\texttt{FILTER (?verdictFitness >= 0.95)} clause applied to actual conformance verdicts in the
ontology. The 2016 lesson---that consequence cannot be recovered from surface form---is encoded
directly into the pipeline's architecture.

\subsection{Enterprise Process Gap}
\label{sec:ggen-lineage-enterprise-gap}

The enterprise process gap is the structural disconnect between the process measurements an
organization produces (event logs, conformance scores, replay traces) and the governance
artifacts it must present to external parties (boards, regulators, acquirers). In the conventional
enterprise, this gap is bridged by human intermediaries who translate technical findings into
business language. The translation is lossy, non-repeatable, and cannot be audited against the
original measurements.

ggen's contribution to closing this gap is the introduction of a \emph{manufacturing contract}:
a declared set of generation rules that specify, for each output artifact, the exact SPARQL query
that extracts its evidentiary basis and the exact Tera template that renders it. The contract is
public, versioned, and machine-executable. Any party with access to the ontology can re-execute
the contract and verify that the manufactured artifact faithfully represents the underlying
measurements.

The board admissibility contract documented in \texttt{README.md} formalizes this through four
pillars: event log integrity and provenance (IEEE 1849-2016 or OCEL 2.0 format, BLAKE3 hash
chain, W3C PROV-O provenance model); mathematical conformance bounds (fitness $\geq 0.95$,
precision $\geq 0.90$); structural model soundness (workflow net soundness under van der Aalst
1998); and cryptographic slide-to-receipt mapping (every slide references a unique BLAKE3 receipt
that can be used to re-run the originating query).

\subsection{Relationship to the Chatman Equation}
\label{sec:ggen-lineage-chatman-equation}

The Chatman Equation, $A = \mu(O^*)$, asserts that admissible artifacts are a function of the
manufacturing process applied to the open ontology. ggen is the most literal instantiation of
this equation in the research foundry. The symbol $\mu$ denotes a manufacturing function; in
ggen's architecture, that function is implemented as the three-layer pipeline: SPARQL extraction
from the ontology, Tera template rendering, and BLAKE3 receipt sealing. The symbol $O^*$ denotes
the open ontology; in ggen's architecture, that ontology is the process intelligence knowledge
base queried through the SPARQL endpoint configured in \texttt{ggen.toml} as
\texttt{endpoint = "in-memory"} with namespaces \texttt{ma:}, \texttt{lifecycle:},
\texttt{wasm4pm:}, and \texttt{compat:}.

The admissibility of the output artifact $A$---whether a board deck slide or a Rust enum
variant---is enforced not by human judgment but by the mathematical properties of the SPARQL
\texttt{FILTER} conditions and the type-law gates in the audit scripts. This is the equation
operating as engineering doctrine: the open ontology is the input, the manufacturing function is
the pipeline, and the admissible artifact is the output.

\subsection{Relationship to Receipts and Replay}
\label{sec:ggen-lineage-receipts-replay}

The receipts-and-replay doctrine holds that every claim must be sealed by a cryptographic receipt
that enables independent verification through replay. ggen implements this doctrine at two levels.

At the \emph{manufacturing level}, each execution of ggen's pipeline is sealed by a BLAKE3 receipt
recording the operation UUID, input hashes, output hashes, and previous receipt hash. The three
receipts in \texttt{.ggen/receipts/} form a hash-chained sequence: the root receipt
(\texttt{sync-20260601-175316.json}, operation \texttt{de2ad19f}) records no predecessor;
the second receipt (\texttt{sync-20260601-203853.json}, operation \texttt{95a9448c}) records
predecessor hash \texttt{612d222\ldots}; the third receipt (\texttt{sync-20260601-230600.json},
operation \texttt{09084e1e}) records predecessor hash \texttt{057521e5\ldots}. This chain makes
the manufacturing history tamper-evident.

At the \emph{artifact level}, every slide in the manufactured M\&A deck references the BLAKE3
receipt hash of the conformance verdict that backs the claim on that slide. A board member who
wishes to verify a synergy projection can take the receipt hash from the slide, locate the
corresponding conformance verdict in the ontology, re-execute the SPARQL query for that verdict,
and confirm that the fitness score and precision score match the values asserted on the slide.
Replay is not a future capability; it is a structural property of the artifact format as declared
in \texttt{ma-deck.tera} and enforced by \texttt{extract-board-claims.rq}.
